\section{Độ trung thực và tính hợp lý}
\label{sec:ch01-do-trung-thuc-va-tinh-hop-ly}

\index{tính hợp lý}%
Dù được đánh giá theo cách nào trong ba cách ở mục trước, một lời giải thích
vẫn có thể thuyết phục được người đọc chỉ vì nó ngắn, quen thuộc và phù hợp
với trực giác. Những phẩm chất đó quan trọng cho giao tiếp, nhưng chưa nói gì
về việc mô hình đã tính toán thế nào. Ta gọi sự phù hợp với đánh giá của
người đọc là \tn{tính hợp lý}{plausibility}; đây là một thuộc tính của lời
giải thích xét trong quan hệ với người nhận.

\index{độ trung thực}%
\begin{definition}[Độ trung thực]
\label{def:ch01-do-trung-thuc}
Lời giải thích có \tn{độ trung thực}{faithfulness} đối với một hành vi và một
tập can thiệp đã xác định khi những đặc trưng mà nó xếp là quan trọng đúng là
những đặc trưng mà hành vi ấy phụ thuộc vào dưới các can thiệp thuộc tập đó,
và những đặc trưng mà nó xếp là không quan trọng thì hành vi ấy không phụ
thuộc vào.
\end{definition}

Chữ phụ thuộc trong định nghĩa~\ref{def:ch01-do-trung-thuc} mang nghĩa của
can thiệp: hành vi phụ thuộc vào một đặc trưng khi một can thiệp được phép lên
đặc trưng ấy làm hành vi thay đổi. Cách phát biểu này là cách sách viết lại
định nghĩa mà phần lớn tài liệu dùng, của Jacovi và Goldberg, theo đó lời
giải thích trung thực là lời giải thích phản ánh đúng quá trình suy luận đằng
sau tiên đoán của mô hình; chương~\ref{ch:chi-so-khong-do-do-trung-thuc} sẽ
dẫn lại định nghĩa ấy qua bài báo neo của cuốn sách~\autocite{p21metrics}.
Sách thêm hai vế hành vi và tập can thiệp vì quá trình suy luận của một mô
hình không quan sát trực tiếp được, còn phản ứng của một hành vi trước một can
thiệp thì đo được.

Định nghĩa~\ref{def:ch01-do-trung-thuc} vì thế buộc ta nêu rõ hai đối tượng
trước khi đánh giá. Thứ nhất, hành vi nào đang được giải thích: một dự đoán, một
lớp đầu ra, hay một vùng đầu vào? Thứ hai, can thiệp nào được phép: che một
đặc trưng, thay nó bằng một giá trị nền, hay thay đổi một khái niệm ở mức
cao? Khi hai công trình trả lời khác nhau cho hai câu này, kết quả của chúng
không nhất thiết mâu thuẫn, vì chúng có thể đã đo hai quan hệ khác nhau.

Tính hợp lý và độ trung thực thường đi cùng nhau trong những ví dụ thuận lợi,
vì vậy chúng dễ bị lẫn; nhưng chúng tách ra ngay khi lời giải thích kể một
câu chuyện quen thuộc mà mô hình không hề dựa vào. Một phương pháp hậu nghiệm
có thể cho lời giải thích mạch lạc trong khi bỏ qua đặc trưng thực sự chi
phối dự đoán. Billa và cộng sự nêu cùng điểm ấy ở dạng hẹp hơn: vì phương pháp
hậu nghiệm chỉ là xấp xỉ, nó có thể không phản ánh đúng cách mô hình đã ra
quyết định~\autocite{p15microscope}.

Từ đây, mỗi lần cuốn sách dùng chữ attribution, câu hỏi đi kèm sẽ là:
attribution đang được xem như một cách trình bày cho người đọc, hay như bằng
chứng về cơ chế? Phần~II của cuốn sách khảo sát các phương pháp trả lời câu
hỏi này bằng những giả định khác nhau.
